% sections/2_theory.tex
\section{Marco Teórico}
El sistema se describe mediante un conjunto de variables de espín $s_i \in \{-1, +1\}$ dispuestas en una red cuadrada bidimensional de lado $L$, con $N = L^2$ sitios. El hamiltoniano del sistema, considerando interacciones a primeros vecinos y la interacción Zeeman con un campo magnético externo $H$, viene dado por:
\begin{equation}
    \mathcal{H} = -J \sum_{\langle i,j \rangle} s_i s_j - \mu H \sum_{i=1}^N s_i
\end{equation}
donde $J>0$ es la constante de acoplamiento ferromagnético y $\mu$ es el momento magnético de cada espín. El parámetro de orden macroscópico, la magnetización total de una configuración, se define como la suma de todos los espines del sistema: 
\begin{equation}
    M = \sum_{i=1}^N s_i
\end{equation}

En el caso particular de ausencia de campo magnético ($H = 0$), la solución analítica exacta obtenida por \citet{PhysRev.65.117} predice una transición de fase continua a una temperatura crítica $T_c$ dada por:
\begin{equation}
    k_B T_c = \frac{2J}{\ln(1+\sqrt{2})} \approx 2.269 J
\end{equation}
adoptando un sistema de unidades naturales donde la constante de Boltzmann es $k_B = 1$. 

En simulaciones de tamaño finito y en ausencia de un campo externo que rompa la simetría, el sistema puede transicionar globalmente entre estados de magnetización positiva y negativa. Para evitar que el promedio térmico se anule artificialmente por debajo de $T_c$, se evalúa el valor esperado del valor absoluto de la magnetización por espín:
\begin{equation}
    m = \frac{\langle |M| \rangle}{N}
\end{equation}

Las fluctuaciones de la energía interna y de la magnetización determinan, a través del teorema de fluctuación-disipación, la capacidad calorífica $C$ y la susceptibilidad magnética $\chi$, respectivamente:
\begin{equation}
    \frac{C}{N} = \frac{\langle E^2 \rangle - \langle E \rangle^2}{N k_B T^2}, \quad \quad \frac{\chi}{N} = \frac{\langle M^2 \rangle - \langle |M| \rangle^2}{N k_B T}
\end{equation}

En la vecindad de $T_c$, la teoría de escalado de tamaño finito predice que estas magnitudes divergen. La capacidad calorífica máxima diverge logarítmicamente, mientras que la susceptibilidad máxima escala según una ley de potencias:
\begin{equation}
\frac{C_{\text{máx}}}{N} \propto \ln L, \qquad 
    \frac{\chi_{\text{máx}}}{N} \propto L^{\gamma/\nu}
    \label{eq:c}
\end{equation}
donde $\nu = 1$ es el exponente asociado a la longitud de correlación y $\gamma = 7/4$ es el exponente de la susceptibilidad magnética. De esta manera, al aplicar logaritmos, esta relación se linealiza para su comprobación empírica como:
\begin{equation}
    \ln\left(\frac{\chi_{\text{máx}}}{N}\right) \propto 1.75 \ln L + \rm cte
    \label{eq:ln_chi}
\end{equation}

Por otro lado, cuando el sistema opera por debajo de la temperatura crítica ($T < T_c$) en presencia de un campo magnético variable, presenta una fuerte dependencia de su historia térmica y magnética debido a la aparición de estados metaestables. Al someter al material a un campo magnético cíclico ($H \rightarrow -H \rightarrow H$), el parámetro de orden no sigue una trayectoria termodinámica reversible, sino que describe un ciclo de histéresis.

Finalmente, la identificación de la transición termodinámica puede formularse matemáticamente como un problema de clasificación binaria. Las variables macroscópicas estocásticas del sistema (energía $E$, magnetización $M$ y susceptibilidad $\chi$) actúan como vectores de características (\textit{features}). Esto permite a un algoritmo de regresión logística estimar la probabilidad posterior de que una configuración pertenezca a la fase desordenada paramagnética ($T > T_c$).